\documentclass{article}
\usepackage{geometry}
\geometry{a4paper, margin=1in}
\usepackage{parskip}
\usepackage{enumitem}
\usepackage{amsmath}
\usepackage{amssymb}
\usepackage{fontspec}
\setmainfont{Times New Roman}

\title{Where the Fault Lands}
\author{}
\date{}

\begin{document}

\maketitle

\section{Causes are not faults}
A cup falls off a table and breaks. Who is at fault?

The wind, through an open window. The person who left the cup near the edge. The person who saw it there and walked past. The person who opened the window. Each of these is a point at which the sequence could have been interrupted, and each is therefore a cause.

None of them is yet a fault. Causes are the things that contributed. Fault is a judgement about which contribution should count against someone, and that judgement is not contained anywhere in the list of causes. Somebody has to make it.

Usually this causes no trouble, because rules and expectations make the judgement for us. Nobody blames the wind, because the wind has no duties. We look for the person who had a duty here, and where the duty is clear the answer is clear, and nobody argues about a cup.

The difficulty appears where no duty has been agreed. No set of agreements can cover everything, and every rule that gets written also marks out the area where nothing has been written. Two people who share a house, a business, or a child produce results together all day, and only a small part of those results falls under anything they have agreed. Where nothing has been agreed, neither the facts nor the rules decide where the fault lands. What decides it is whoever happens to be doing the judging.

Consider a kind of case that occurs in hospitals: a nurse gives a patient the wrong drug, because two medicines are supplied in almost identical containers. The nurse administered it. The pharmacy stocked both. The manufacturer designed the packaging. The hospital set the staffing level for that shift. Each is a point where the sequence could have been broken, and for a long time the answer given was the nurse. Aviation and medicine have both increasingly treated a single named cause as a sign that the investigation stopped too early---not out of fairness, but because blaming individuals produced unreliable information about what had actually happened.

\section{How to judge well}
If neither the facts nor the rules decide it, then working out who is at fault is a judgement made without enough information. That is not a reason to stop making such judgements. It is a reason to make them the way careful people make any judgement of that kind: weigh the evidence, keep the conclusion in proportion to it, change it when new information appears, and let other people examine the result.

Two of these requirements matter more than the others.

The first is that you must keep yourself among the possible causes. If your actions helped produce the result, you are one of the candidates. Removing yourself before looking is not modesty. It is a search that cannot succeed, because the right answer may already have been taken out of it, and no amount of careful thinking afterwards will recover it. Everything that follows will be a conclusion about the remaining candidates only.

The second is that you must be able to leave a question unresolved. Often there is not enough information to decide, and the honest position is that it is not yet decided. Forcing a conclusion at that point does not produce a well-supported answer. It produces an answer supplied by something other than the evidence: by habit, by what you expected, or by what suits you.

Both requirements are normally treated as matters of character, as being fair or being willing to admit things. They are better described as matters of competence. Someone who removes themselves from the list before looking, and who cannot tolerate an unresolved question, is not being unkind. They are conducting a search that cannot reach the right answer, and then reporting the result with confidence.

\section{When one answer is unavailable}
Some people cannot afford to be at fault. Three quite different situations produce this.

For some, being at fault does not stay attached to what they did. Ordinary guilt has an ending: you acknowledge the thing, you repair it, and it is finished, because it was always about the act. Where it attaches to the person instead, nothing finishes it, because nothing repairs being something. The sign of this is that size stops mattering. Admitting a small thing costs as much as admitting a large one, which is why such a person will resist most fiercely over matters that look trivial to everyone watching.

For others, the expense comes from their position. A professional whose licence is at risk. A company facing a claim. Anyone whose account of events will later be used against them.

For others again, the cost is simply what they stand to lose. Where the present arrangement suits someone, conceding a point puts it at risk.

These three differ in nearly every respect: in what the person feels, in what they are aware of, in what alternatives they have. They share one feature, and it is the only one that matters here. In each case an answer has become unavailable before the evidence is considered.

\section{What follows}
None of what follows requires deliberate invention. People do knowingly construct excuses, and that is a real thing, but it is not the interesting case and it is not the durable one. If a search begins with one candidate already removed, ordinary evidence is usually enough to support whatever remains. A man who cannot have left the gate open will remember that his son was out that afternoon. This will be true, and it will occur to him as a discovery rather than as a defence. The consequence is that his sincerity tells you nothing at all: a search narrowed in advance produces conclusions that feel exactly like conclusions reached properly, because in every respect other than the narrowing, they were.

Blame then moves from the event to the person. Saying that someone was careless on Tuesday settles nothing, because the same question opens again the following week. Saying that someone is careless settles every instance in advance, including those that have not happened yet. An explanation that has to survive repetition will therefore describe the person rather than the occasion, and a description of a person covers far more ground than any single event. What gets moved is always larger than what caused it to move. This is not a decision to be harsh. It is what a durable explanation has to look like.

Changing your mind then becomes more expensive each time. The explanations accumulate and they have to agree with one another, so withdrawing one puts the others in question. The cost of correcting course rises with every repetition, until continuing costs less than stopping, and after that the gap only widens.

The account also has to be maintained. An explanation that has not been confirmed for some time begins to look like an accusation rather than a fact, and an accusation can be disputed. Confirming instances are therefore needed, and where none occur, the definition of what counts as a fault widens until some appear. This gives a prediction that can be checked: a period in which the other person does nothing wrong will not produce fewer complaints. It will produce complaints about smaller things.

Finally, a description of someone's character becomes grounds for acting on it, and the steps are worth setting out because each follows from the one before. You did this. You are the kind of person who does this. It will therefore happen again. I must protect against it. Restricting you is justified. And then the step that closes the circle: the restriction produces friction, the friction produces incidents, and the incidents confirm the description that justified the restriction in the first place.

None of this is confined to households. A public body discovers that something under its supervision has gone badly wrong. It must produce an account, and its own procedures are among the possible causes, but that answer carries consequences it cannot absorb, so the account settles on a contractor. The first version describes what the contractor did. Later versions describe what the contractor is like, because that version also covers the next incident. Each subsequent report is read in the light of the earlier ones, and each makes withdrawing the original finding more costly. Eventually the body restricts the contractor, which is by then not merely permitted but required, since the body has itself established the risk. The restrictions make the contractor's work harder, and the resulting failures are recorded. No one in this sequence has to act in bad faith for it to run.

\section{Why nothing corrects it}
Most mistaken blame is harmless, because somebody disagrees with it. Being wrong about who is at fault is ordinary, and what makes it survivable is that other people are watching and will say so. Correction does not come from inside a person. It arrives from outside.

Here it does not arrive, and the person being blamed cannot supply it. A description broad enough to cover a person rather than an occasion can absorb almost any response to it. Denial is what the description predicts. Silence is taken for agreement. Accepting small things to keep the peace supplies exactly the instances the account requires. The difficulty is not that such a description must be false. It is that the process has stopped producing observations that could show it to be false.

There is a further difficulty, which concerns what evidence is available at all. Being the only person who is told something is a considerable advantage in any dispute about fault. Information that does not reach someone cannot be acted on by them, and their failure to act is then real and can be described accurately. Shared acquaintances who hear one account and not the other stop being witnesses and become an audience. Outside sources, such as a professional opinion or a written agreement, can be set aside before they are consulted, by describing the source rather than disputing what it says. None of this requires a single false statement. It requires only that the supply of evidence narrows.

The person who keeps to the two requirements is therefore in the worst position of all. Staying among the possible causes means fault can land on them, so it does, and each instance becomes a precedent for the next. Meanwhile the evidence they are weighing so carefully has been produced by the process they are standing inside. Careful reasoning applied to a controlled supply of evidence produces a confident wrong answer. Fairness is not merely unrewarded in these conditions. It is a resource the process consumes, and being fair-minded can be exploited wherever somebody else decides what you are shown.

How it ends depends on how difficult it is to leave. Where leaving is possible, this usually ends it, though the cost of leaving is rarely only a formal matter. Where leaving is not possible, because of a child, a job, or a shared border, it does not end, and instead the second person's room to act keeps narrowing while what is expected of them stays the same.

\section{How to tell the difference}
If neither the facts nor the rules decide who is at fault, how can anyone distinguish a considered conclusion from a fixed one?

Watch what the conclusion does when something arrives that does not fit it. A conclusion that follows the evidence can take an exception without much trouble. It narrows, or it acquires a qualification, or it is set aside for a better one. A conclusion fixed before the evidence has to do something else. It widens, so that the exception is covered as well. It picks up additional assumptions whose only function is to protect it. And it becomes more costly to withdraw with each repetition, because withdrawing any part of it now puts the rest in doubt. The signature is not confidence, and it is not cost by itself. It is that the conclusion never becomes easier to abandon.

This can only be seen from outside. It requires a long record, and some view of what the conclusion is costing the person who holds it, and nobody inside the situation has either. Correction can arrive by other routes---through documentation, a formal procedure, or consequences that fall where the account did not predict---but each of these depends on some part of the evidence escaping the control of the people producing it. Where the participants cannot correct the evidence themselves, observation from outside is what the whole thing rests on.

None of this is specific to people. It follows from two facts: that those who act together produce results neither produced alone, and that deciding who is responsible for such results is a judgement rather than an observation. Any system of agents in that position meets the same problem, whether individuals, companies, governments, or artificial systems that reason and then act on their conclusions. Wherever being at fault becomes expensive enough, reasoning about fault stops working, and it stops working in a way that feels from the inside like being careful.

\section{Conclusion}
A group does not stay correctable because its members are willing to be fair. It stays correctable only while responsibility can still be assigned to anyone within it, including whoever is doing the assigning. Not because everyone is equally responsible for everything, but because the possibility has to remain open until the evidence has been looked at. Where it closes for one participant, it closes for the process.

\end{document}